% Lösungen Einheit 1 von 3 – Figur und Term (zusammenbau v0.1)
\einheitenkopf{Einheit 1 von 3 · Figur und Term}

\erg{9}{a) $f(u)$ – die Höhe ist der Funktionswert an der Stelle $u$ \quad b) Breite $2$ (die Stelle), Höhe $f(2) = 5$ (der Funktionswert) \quad c) Ecken $(-2 | -3)$, $(2 | -3)$, $(2 | 3)$, $(-2 | 3)$; Seiten $2 \cdot 2 = 4$ cm und $2 \cdot f(2) = 6$ cm, Fläche $4 \cdot 6 = 24$ cm² \quad d) Rechter Winkel bei $P$: Grundseite $a$, Höhe $f(a) = a^2 \cdot e^{-a}$, also $A = \tfrac{1}{2} \cdot a \cdot a^2 \cdot e^{-a} = \tfrac{1}{2} a^3 \cdot e^{-a}$. \quad e) Parallele Seiten $\overline{AB} = 10$ und $\overline{D_uC_u} = 10 - 2u$, Mittelparallele $\frac{10 + 10 - 2u}{2} = 10 - u$, Höhe $f(u)$ – Fläche ist Mittelparallele mal Höhe. Schenkel: $\sqrt{u^2 + (f(u))^2}$ (Pythagoras mit Überstand $u$ und Höhe $f(u)$). \quad f) Für $k = 3$: $A_k(-2 | 3)$, $B_k(2 | 3)$, $C(0 | 5)$. Für $k$ gegen 1 schrumpft die Grundseite, für $k$ gegen 5 die Höhe gegen null; beide Ränder sind ausgeschlossen, also gibt es kein kleinstes Dreieck, dazwischen aber ein größtes. $A(u) = \tfrac{1}{2} \cdot 2u \cdot (5 - f(u)) = 4u - 0{,}5u^3$}
\erg{10}{a) Die Höhe ist der Funktionswert $f(a) = 6 - a$, nicht die Stelle $a$: $A = \tfrac{1}{2} \cdot a \cdot (6 - a)$. \quad b) Die waagerechte Seite reicht auf der x-Achse von 0 bis zur Stelle $u$, die senkrechte von der x-Achse bis zur Höhe des Graphenpunkts, also bis $f(u)$. \quad c) $P(2 | 4)$, Fläche $2 \cdot 4 = 8$; allgemein $A(u) = u \cdot (8 - u^2) = 8u - u^3$}
